\section{Introduction}\label{sec:intro}

The Yang--Mills existence and mass gap problem asks for a quantum
Yang--Mills theory on $\R^4$ with a compact simple gauge group that satisfies
the axioms of constructive quantum field theory, and for a number $\Delta>0$
such that the spectrum of the Hamiltonian lies in $\{0\}\cup[\Delta,\infty)$
\cite{JW06}. In the Euclidean formulation the theory is a probability measure
on gauge fields whose Schwinger functions satisfy the Osterwalder--Schrader
axioms \cite{OS73,OS75}; the mass gap is then exponential decay of the
truncated two-point functions of gauge-invariant observables at the rate
$\Delta$. The natural regularisation is Wilson's lattice gauge theory
\cite{Wil74}: an element $Q_e$ of $G\subset\U(N)$ sits on every link $e$ of the
lattice $a\Z^4$, and a configuration has the weight
$\exp(N\beta\sum_p\ReTr Q_p)$, where $Q_p$ is the ordered product around the
plaquette~$p$.\footnote{This is the normalisation of~\cite{SZZ23,BB26a}. The
conventional inverse coupling is $\beta_W=N^2\beta$, and for $\SU(2)$,
identified with the unit quaternions, $\beta_W=4\beta$.} Asymptotic freedom
forces the continuum limit $a\to0$ to be taken with $\beta\to\infty$, that is,
at weak coupling.

At small $\beta$ the lattice theory has a unique state, exponential clustering
and a transfer-matrix gap, by convergent expansions \cite{OSei78,Sei82};
Ba{\l}aban's renormalisation group \cite{Bal84a,Bal85a,Bal89b} gives ultraviolet
stability in finite volume at weak coupling (Dimock \cite{Dim13} presents the
method for the scalar $\phi^4$ model in three dimensions), and Magnen, Rivasseau and S\'en\'eor constructed the theory with an
infrared cutoff \cite{MRS93}. No existing result gives a mass gap in physical
units. This paper proves a lattice mass gap with explicit constants for every
compact simple group, reduces the continuum gap to one estimate on one
physical scale, proves that estimate at strong coupling, reduces it further to
a decay of correlations across a slab of physical width $\ell/2$, and
determines what the methods used here can and cannot reach. It also states
precisely what is missing.

\subsection{The lattice gap with explicit constants}

Let $G$ be a compact connected Lie group with semisimple Lie algebra,
identified with its image under a faithful unitary representation in $\U(N)$,
and let $\kappa_G>0$ be its least Ricci curvature for the bi-invariant metric
$\ip XY=-\Tr(XY)$. For a simple group $\kappa_G=h^\vee/(2\ell_\rho)$, with $h^\vee$
the dual Coxeter number and $\ell_\rho$ the Dynkin index of the representation;
for $\SU(N)$, $\kappa_G=N/2$.

\begin{mainthm}\label{thm:A}
Suppose $N|\beta|<\frac2{29}\kappa_G$. For every gauge-invariant local
observable the expectations in boxes with free boundary converge, along every
exhausting sequence of boxes, to a state $\omega$ that is invariant under the
lattice symmetries; truncated correlations of Lipschitz observables decay at
every rate below
\[
m_G(\beta)=\operatorname{arccosh}\Bigl(1+\frac1{12}\Bigl(\frac{\kappa_G}{N|\beta|}
-\frac{29}2\Bigr)\Bigr)>0 ,
\]
and for $\beta\ge0$ the state is reflection positive and its transfer matrix
has spectral gap $m\ge m_G(\beta)$. For $G=\SU(2)$ the same holds for
$|\beta_W|<\frac13$ with $m\ge-\log(3\beta_W)$.
\end{mainthm}

Two ingredients carry the proof. The first is a gauge: a field of straight
lines on $\Z^4$ that contains exactly one link of every plaquette
(Sec.~\ref{sec:linefield}). Setting these links to the identity loses no
gauge-invariant information, and every remaining link then interacts with
exactly twelve others instead of eighteen. The second is a Hessian bound in
this gauge, $|\Hess S|\le\frac{29}2$ for the action $S=\sum_p\ReTr Q_p$, proved
by a divergence identity and a positivity statement for one plaquette
(Sec.~\ref{sec:curvature}); with the Ricci curvature of $G$ it gives a
Bakry--\'Emery bound \cite{BE85,BGL14} with an explicit decay rate. Without
the gauge the sharp Hessian bound is $16$ \cite{BB26a}, so the gauge enlarges
the range of every group by the factor $32/29$ and supplies the rate. The
$\SU(2)$ statement combines the gauge with the sharp Kantorovich modulus
$\frac14$ of the von Mises--Fisher family \cite{BB26a} in Dobrushin's
criterion \cite{Dob68,Dob70}.

\subsection{A continuum mass gap from one cube}

The lattice gap of Theorem~\ref{thm:A} is measured in lattice units and
vanishes in physical units as $a\to0$. The second result isolates a single
estimate, posed on one scale, that controls the gap in physical units. On a
periodic lattice let $\mathcal P_x$ resample the links of the cube of side $n$
at $x$ from their conditional law under the Wilson measure $\mu$. The
condition $\BP(n,\varpi)$ asks that
\[
\varpi\Var_\mu(f)\le\frac1{n^4}\sum_x\mu\bigl((f-\mathcal P_xf)^2\bigr)
\]
for every gauge-invariant $f$: the dynamics that resamples cubes, each link
at rate one, has spectral gap at least $\varpi$ on gauge-invariant functions.
The condition involves no gauge fixing and no supremum over boundary
conditions.

\begin{mainthm}\label{thm:B}
Let $a_k\to0$ and $\beta_k\ge0$ be couplings, and suppose that there are
$\varpi>0$, $\ell>0$ and integers $n_k\ge2$ with $n_ka_k\le\ell$ such that
$\BP(n_k,\varpi)$ holds for the Wilson measure at $\beta_k$ on all large tori.
Then every continuum limit of the corresponding lattice theories has
Hamiltonian $H$ with $\sigma(H)\subset\{0\}\cup[\Delta,\infty)$,
\[
\Delta=\frac{16\varpi}{625e^2\ell},
\]
and $0$ is a simple eigenvalue; if the limit is nontrivial, $H$ has a mass
gap of at least $\Delta$.
\end{mainthm}

The proof (Sec.~\ref{sec:cube}) turns the spectral gap into decay of
covariances through a finite-speed bound for the cube dynamics, at a rate
of order $1/n$ in lattice units, hence of order $1/\ell$ in physical units,
and passes this decay through the Osterwalder--Schrader reconstruction.

\subsection{Where the cube condition holds}

\begin{mainthm}\label{thm:C}
\begin{enumerate}[label=\textup{(\roman*)},leftmargin=*]
\item At $\beta=0$, $\BP(n,1)$ holds for every $n$. For
$21N|\beta|<\kappa_G$ and $n$ large, $\BP(n,1-\alpha_n)$ holds with
$\alpha_n=N|\beta|(18n-3)/((\kappa_G-21N|\beta|)n^2)\to0$; for $\SU(2)$ the
range is $|\beta_W|<\frac5{18}$, and for $\SU(3)$ it contains
$0\le\beta_W\le0.45$.
\item For the free scalar field of lattice mass $\mu$ in $d$ dimensions the
best constant is
$\varpi=\frac{\mu^2}{n^d}\ip{\mathbf 1}{(-\Delta_D+\mu^2)^{-1}\mathbf 1}$, with
$-\Delta_D$ the Dirichlet Laplacian of the cube, and with $s=n\mu$ (the mass
times the physical size of the cube)
\[
\Bigl(\frac56\Bigr)^d\frac{s^2}{s^2+10d}\le\varpi\le\frac{s^2}{s^2+16d/9}.
\]
\item For $G=\SU(2)$ and $|\beta_W|\le\frac7{20}$, $\BP(n,\varpi)$ holds for
every $n\ge2$ with $\varpi=\frac{106}{1933}\cdot517^{-4}$, independent of $n$
and of the torus, and $\BP(n,1-259/n)$ holds for $n\ge260$. For
$0\le\beta_W\le\frac9{25}$ every torus limit is reflection positive and its
transfer matrix has spectral gap at least $7/(625\cdot2957\,e^2)$.
\end{enumerate}
\end{mainthm}

Part (ii) is the model case of the continuum tuning: the constant is uniform
in the lattice spacing at fixed $m\ell$, and it vanishes as $m\to0$. The cube
condition is therefore a quantitative statement of mass generation at the
scale $\ell$. Part (i) is proved in Sec.~\ref{sec:bpstrong} by the gauge of
Theorem~\ref{thm:A} and a weighted Dobrushin criterion for blocks of links.
No argument built on single-link updates in this gauge passes
$\beta_W=\frac13$ for $\SU(2)$ (Proposition~\ref{prop:methodlimit}). Part
(iii) passes that value by resampling \emph{stars}, the eight links at a
vertex of one parity (Sec.~\ref{sec:stars}): the perturbation of one
plaquette at a vertex enters the law of a star only through the relative
variable of two of its links, and no plaquette is internal to a star.

\subsection{Mass generation at one scale}

Statement (D) below asks for the cube condition along the continuum tuning.
Section~\ref{sec:massgen} reduces it to a hypothesis of a different kind,
closer to physics. Split a box of the lattice, all of whose widths lie
between $2d-1$ and $md-1$, into two outer parts of width $d$ and a middle of
width $d-1$ in one coordinate, fix an arbitrary configuration outside the
box, and let $\rho$ be the maximal correlation between the gauge-invariant
observables of the two outer parts under the conditional law. The hypothesis
$(\mathrm M_{d,m}(c))$ asks that $\rho\le c$ for all such boxes, splittings
and boundary conditions. Along the tuning $d\approx\ell/(2a)$, so this is a
decay of correlations across a slab of physical width about $\ell/2$,
uniform in the boundary condition.

Theorem~\ref{thm:D} proves that $(\mathrm M_{d,m}(c))$ with $c<\frac12$ and $m$
large in terms of $c$ gives $\BP(2d-1,(1-2c)^4/32)$, hence, along the tuning,
statement (D) and the continuum gap $(1-2c)^4/(1250\,e^2\ell)$ in every
nontrivial limit. The proof passes through a gap for the dynamics that
resamples overlapping patches of a box, obtained from the decorrelation of
the slabs by a Markov factorisation, and a finite-size criterion for
gauge-invariant functions in the spirit of those of Knabe~\cite{Kna88} and of
Gosset and Mozgunov~\cite{GM16} for quantum spin systems. The hypothesis also
gives exponential decay of gauge-invariant correlations under every boundary
condition, at rate $\log(1/c)/d$ per lattice spacing
(Proposition~\ref{prop:alldecay}). It holds at strong coupling at fixed $d$ and
$m$, and in two dimensions along the weak-coupling tuning
(Theorem~\ref{thm:twod}). In four dimensions it is not proved.

\subsection{The limits of the methods}

Appendix~\ref{app:obstructions} proves that several natural routes to (D)
cannot work, and locates the threshold of Theorem~\ref{thm:D}.
\begin{enumerate}[label=\textup{(\alph*)},leftmargin=*]
\item \emph{The threshold $\frac12$ is sharp.} For a massless Gaussian chain
the maximal correlation of the slab mechanism equals $\frac12$ exactly, and
the gap of the corresponding block dynamics on $M$ blocks is at most
$16/(3M)$; so no decorrelation constant $c\ge\frac12$ can give a gap uniform in
the size (Proposition~\ref{prop:gausschain}).
\item \emph{Two dimensions.} For $\SU(2)$ in two dimensions the middle slab of
a box integrates out, after an axial gauge, to a convolution power of the
heat kernel, and $(\mathrm M_{d,m}(c))$ holds with
$c=20\,e^{-3g^2(da)^2/4}$, which tends to zero as the physical size of the box
grows (Theorem~\ref{thm:twod}). The same condition for strips one link wide
fails along the tuning: a pinned strip is a Brownian bridge on $\SU(2)$, and
the antipodal boundary condition breaks a global symmetry, with correlation
close to one (Proposition~\ref{prop:strip}).
\item \emph{Curvature.} The Bakry--\'Emery constant of the Wilson measure
without gauge fixing is at most $N/2-16N|\beta|$, so the Hessian bound $16$ of
\cite{BB26a} is sharp (Theorem~\ref{thm:fullspace}); in the gauge of
Theorem~\ref{thm:A} the constant $\frac{29}2$ cannot be replaced by anything
below $\frac{38}3$, and no Bakry--\'Emery argument for the gauge-fixed measure
reaches $|\beta_W|=\frac3{19}$ (Theorem~\ref{thm:slice}). Every smooth
plaquette action, such as an effective action at a coarse scale, obeys a bound
of the same kind (Proposition~\ref{prop:coarse}).
\item \emph{Entropy, covariances and fixed blocks.} Block entropy
factorisation implies the cube condition and so cannot be easier
(Proposition~\ref{prop:entbp}); the link covariances of every gauge-invariant
measure are $\frac14I$ and carry no information on the coupling
(Proposition~\ref{prop:linkcov}); and a single-link spectral gap, or any gap
for blocks of fixed lattice size, with constant $C_k$ satisfying $C_ka_k\to0$,
forces every continuum limit to be trivial
(Proposition~\ref{prop:fixedblocks}).
\item \emph{A gap in the limit is not enough.} If the lattice gaps $m_k$ tend
to zero, a choice of lattice spacings proportional to $m_k$ produces a
nontrivial continuum limit with a mass gap; this limit is kinematic and is not
identified with Yang--Mills theory (Theorem~\ref{thm:gaptuning}). The content
of (D) and (E) is that the gap and the nontriviality hold along the tuning of
asymptotic freedom.
\end{enumerate}

\subsection{What remains}

Theorems~\ref{thm:A} and~\ref{thm:C}(i), (iii) are strong-coupling theorems.
The Clay problem needs the opposite end of the coupling range: the continuum
limit lives at $\beta\to\infty$, which the ranges of these theorems do not
reach. Two
statements would complete the proof of a continuum mass gap, and neither is
proved here:
\begin{enumerate}[label=\textup{(\Alph*)},leftmargin=*]
\setcounter{enumi}{3}
\item $\BP(n_k,\varpi)$ with $n_ka_k\le\ell$ along the asymptotic-freedom
tuning $\beta_k\to\infty$, uniformly in $k$;
\item convergence of the renormalised Schwinger functions to a nontrivial
limit satisfying the Osterwalder--Schrader axioms.
\end{enumerate}
Given (D) and (E), Theorem~\ref{thm:B} gives the mass gap $\Delta\ge
16\varpi/(625e^2\ell)$, and by Theorem~\ref{thm:D} statement (D) follows from
$(\mathrm M_{d_k,m}(c))$ along the tuning with $c<\frac12$. Section~\ref{sec:weak} explains why (D) is the heart of
the matter: by Theorem~\ref{thm:C}(ii) it cannot hold without a dynamically
generated mass, it fails at fixed cube size at weak coupling, and it fails in
the Coulomb phase of the abelian theory, so its proof has to use the
non-abelian structure of $G$ at every scale up to $\ell$.

\subsection{Related work}

Strong-coupling clustering goes back to Osterwalder and Seiler \cite{OSei78}
and Seiler \cite{Sei82}; probabilistic treatments include the survey of
Chatterjee \cite{Chat19}, the Langevin approach of Shen, Zhu and Zhu
\cite{SZZ23}, whose Bakry--\'Emery argument we sharpen, the area law in
enlarged regions of Cao, Nissim and Sheffield \cite{CNS25}, and correlation
decay for finite gauge groups at weak coupling by Adhikari and Cao \cite{AC25}.
In two dimensions the continuum measure is known
\cite{GKS89,Dri89,Lev03}. The Kantorovich modulus of the $\SU(2)$ single-link
law and the sharp Hessian bound without gauge fixing are proved in
\cite{BB26a}, and an earlier and much longer account of the one-cube
reduction is \cite{BBLR26}.

One of us (I.~J.~S.) has formalised in Lean~4 with Mathlib~\cite{Ses26}, without
unproved statements and with only the standard axioms, several ingredients of
the strong-coupling analysis on periodic lattices: the Haar probability on
$\SU(N)$, an infinite-volume Gibbs state obtained as a limit of torus states,
site reflection positivity and the transfer matrix, the Holley--Stroock
comparison, Dobrushin's covariance bounds, and the passage from decay of
correlations to a transfer-matrix gap (one direction of
Proposition~\ref{prop:gapdecay}). For the periodic state it proves a
transfer-matrix gap for every $N$ at $\beta_W<5.9\cdot10^{-5}$, and explicit
single-link heat-bath gaps on boxes of side fifteen for $\SU(2)$ at
$\beta_W\le0.166$ and for $\SU(3)$ at $\beta_W\le0.38$ (the coupling there is
$\beta_W/2$). It also proves, by a Poincar\'e inequality and a light-cone
argument different from Lemma~\ref{lem:speed}, that an ungauged one-cube
influence condition gives a lattice gap of order $1/n$; this is a variant of
the lattice step, Theorem~\ref{thm:decay} and Corollary~\ref{cor:latticegap},
at a smaller rate. The continuum step, the gauge
fixing and the curvature bounds are not part of that development, which shares
with this paper the open hypothesis (D).

